% Generated draft: requires mathematical and source review.
\documentclass[11pt]{article}
\usepackage[margin=25mm]{geometry}
\usepackage{fontspec}
\setmainfont{Latin Modern Roman}
\usepackage{amsmath,amssymb,mathtools,unicode-math}
\setmathfont{Latin Modern Math}
\usepackage{enumitem,needspace,xurl,hyperref}
\setlength{\parindent}{0pt}
\setlength{\parskip}{5pt}
\setlength{\emergencystretch}{3em}
\newcommand{\N}{\mathbb{N}}
\newcommand{\Z}{\mathbb{Z}}
\newcommand{\Q}{\mathbb{Q}}
\newcommand{\R}{\mathbb{R}}
\newcommand{\C}{\mathbb{C}}
\newcommand{\F}{\mathbb{F}}
\newcommand{\A}{\mathbb{A}}
\newcommand{\PP}{\mathbb{P}}
\newcommand{\E}{\mathbb{E}}
\newcommand{\eps}{\varepsilon}
\newcommand{\dd}{\,\mathrm{d}}
\begin{document}
\section*{UnsolvedMath ID: 30001738}
\textbf{Source problem code:} catalogue entry
\subsection{Definitions and mathematical statement}
Let \(E/F\) be a quadratic extension of local fields, with nontrivial Galois involution \(\tau\). For a positive integer \(n\), put \(G=\operatorname{GL}_n(E)\) and let
\[
X=\{x\in\operatorname{GL}_n(E):{}^t\tau(x)=x\}
\]
be the set of nonsingular Hermitian matrices, where \(\tau\) acts entrywise. Its natural right action is
\[
x\cdot g={}^t\tau(g)\,xg\qquad(x\in X,\ g\in G).
\]
Write \(n=n_1+\cdots+n_k\), with \(k\geq1\) and \(n_i\geq1\), so that \(\delta_i\) is a representation of \(\operatorname{GL}_{n_i}(E)\). The notation
\[
\pi=\delta_1\times\cdots\times\delta_k
\]
denotes parabolic induction from these blocks; normalized parabolic induction is the conventional interpretation, but the supplied record does not explicitly specify its normalization. Each \(\delta_i\) is assumed to be essentially square-integrable, meaning that a twist by a character is square-integrable modulo the center, in the intended representation category. Define \(\pi^\tau(g)=\pi(\tau(g))\); \(\tau\)-invariance means \(\pi^\tau\simeq\pi\).

\textbf{Source conjecture.} For every such product \(\pi\) that is \(\tau\)-invariant, is
\[
\dim\operatorname{Hom}_G(X,\pi^*)=2^k?
\]
Here \(\pi^*\) is provisionally read as the contragredient representation. The hypothesis concerns invariance of the entire induced representation; it does not assert invariance of each factor.

\textbf{Unresolved qualifications.} The displayed formula retains the source notation. Since \(X\) is a variety or a set rather than a representation, \(\operatorname{Hom}_G(X,\pi^*)\) requires a specified function space, distribution space, or other representation associated with \(X\), together with an equivariance and continuity convention. None is supplied, so the multiplicity space cannot yet be defined unambiguously. The record also leaves unspecified whether the local fields are Archimedean or non-Archimedean, any characteristic restrictions, the representation category and coefficient field, and the precise dual denoted by \(\pi^*\). It does not explicitly require \(\pi\) to be irreducible or impose an ordering or additional compatibility conditions on the inducing factors. These missing qualifications must be checked against the report before this becomes a fully precise statement; no additional hypotheses are imposed here.
\subsection{Short English statement}
The supplied conjecture predicts that a Galois-invariant representation induced from \(k\) essentially square-integrable factors has multiplicity \(2^k\) for the space of nonsingular Hermitian matrices. The representation associated with that space and several representation-theoretic conventions remain unspecified. This is a statement-only draft requiring human mathematical review.
\subsection{Sources}
\begin{itemize}
\item Supplied local UnsolvedMath registry record, permanent ID \(30001738\), title: Multiplicity Formulas for Galois-Invariant Induced Representations; exact source problem code: \texttt{catalogue entry}. The record belongs to the Oberwolfach Reports Open Problems collection and gives proposed year \(2011\); no individual proposer is identified.
\item \emph{Automorphic Forms: New Directions} (2011), Oberwolfach Reports, DOI \texttt{10.4171/owr/2011/14}, pp.\ 719--779. Supplied source URL and literature-source link, labeled ``original OWR report'': \url{https://doi.org/10.4171/owr/2011/14}. The supplied background describes Oberwolfach Reports as records of mathematical workshops and their open-problem discussions. No specific contribution, author, or problem page within this report is supplied.
\item Literature assessment dated \texttt{2026-08-21}: the supplied curation labels the problem open, describes the general multiplicity formula as conjectural, and states that the cited report proves only special classes, such as ladders. This is supplied dated curation, not an independently verified literature assessment; its status and cited claims require independent checking.
\item The registry labels the statement \texttt{corrected\_verified} and reports that its mathematical content was checked against the official report, with notation, context, and question form normalized without changing the source-level claim. This verification assertion is supplied metadata; the present draft does not independently verify it or claim to have consulted the report.
\end{itemize}
\end{document}
